\chapter{Aide-mémoire clavier et souris}
\label{chap:raccourcis}
Les gestes sont également expliqués dans leurs familles. Leur effet dépend
du focus : une flèche dans un curseur ajuste sa valeur ; ce n’est pas une
commande globale de changement de photo. Les raccourcis ci-dessous reprennent
les séquences déclarées dans la version décrite, notamment sur Windows.

\begin{longtable}{@{}L{4.4cm}L{9.2cm}@{}}
\toprule Touche ou geste & Action et contexte\\\midrule\endhead
\touche{Ctrl + virgule} & Ouvrir les préférences.\\
\touche{Ctrl + +} & Agrandir une photo prête. \touche{Ctrl + =} est également accepté.\\
\touche{Ctrl + $-$} & Réduire une photo prête.\\
\touche{Ctrl + 0} & Ajuster à la vue.\\
\touche{Ctrl + 1} & Afficher à 100\,\% de l’aperçu.\\
\touche{Ctrl + Maj + C} & Activer ou désactiver le comparatif partagé, si les deux aperçus sont disponibles.\\
\touche{Espace} maintenu & Original temporaire lorsque la vue image a le focus ; relâcher restaure la vue.\\
\touche{Ctrl + molette} & Zoom centré autour du pointeur dans la vue image.\\
Molette sans Ctrl & Déplacer l’image si elle dépasse la vue.\\
\touche{Maj + molette} & Déplacement horizontal à partir de la molette verticale, sur une image agrandie.\\
Clic dans la photo & Basculer entre ajustement et agrandissement, sans glissement.\\
Clic gauche maintenu, déplacement & Parcourir une image agrandie ; ou déplacer le séparateur du comparatif.\\
\bottomrule
\end{longtable}

\clearpage
\section*{Réglages et manipulation des contrôles}
\begin{longtable}{@{}L{4.4cm}L{9.2cm}@{}}
\toprule Touche ou geste & Action et contexte\\\midrule\endhead
\touche{Tab}, \touche{Maj + Tab} & Parcourir les contrôles qui acceptent le focus.\\
\touche{Espace} sur un bouton ou une case & Activer la commande focalisée ; ne pas confondre avec le maintien dans la vue image.\\
Flèches sur un slider & Ajuster au pas déclaré : par exemple 0,01 EV pour l’exposition.\\
\touche{Entrée} dans un champ & Confirmer la saisie numérique.\\
Double-clic sur nom ou slider & Restaurer la référence du paramètre ; cette référence n’est pas toujours zéro.\\
Flèches sur une poignée de panneau & Redimensionner au pas de 10 px dans sa direction.\\
Double-clic sur une poignée de panneau & Restaurer sa dimension usuelle.\\
Double-clic dans les courbes & Ajouter un point de contrôle.\\
Clic droit sur un point intérieur & Supprimer ce point de courbe.\\
Glisser dans une roue de grading & Modifier teinte et intensité de la zone active.\\
\bottomrule
\end{longtable}

\attention{Ne pas supposer les raccourcis de Lightroom ou d’un autre éditeur :
Ctrl + O, Ctrl + S, Ctrl + E, Ctrl + Z et les touches numériques de notation
ne sont pas déclarés comme commandes globales de l’application décrite.
Les raccourcis propres aux dialogues natifs dépendent du système.}
